\section{Summary of Findings}

The experiments demonstrate that fine-tuning improves the performance of both
evaluated model families in Java-to-C++ code translation. For both StarCoder2
and Qwen, the fine-tuned variants achieved higher BLEU scores and a larger
number of accepted solutions compared to the corresponding base models.

Among the evaluated systems, the fine-tuned Qwen model achieved the strongest
overall performance. It obtained the highest average BLEU score and the largest
number of accepted translations in the execution-based evaluation. The
fine-tuned StarCoder2 model also showed considerable improvement over its base
variant, particularly in terms of reducing compilation errors and increasing the
number of executable solutions.

The results also highlight the importance of execution-based evaluation. While
BLEU score was useful for measuring textual similarity between generated and
reference code, it did not always reflect whether the translated programs were
functionally correct. Several outputs with relatively high BLEU scores still
failed to compile or produced incorrect results during testing. This finding is
consistent with previous studies, which report that similarity-based metrics do
not reliably capture semantic correctness in code-related tasks.

Overall, the experiments indicate that fine-tuning can substantially improve
both the syntactic quality and the functional correctness of translated code.
However, the remaining compilation errors, runtime errors, and wrong answers
show that current code translation models are still far from perfect. The
results suggest that future research should focus not only on improving model
architectures, but also on developing stronger evaluation methods that better
reflect real-world correctness and robustness.